% =============================================================================
% 7. Limitations and Future Work -- target ~1.5 pp
% Be honest. Limitations sections that read as marketing read as defensive.
% =============================================================================

\section{Limitations and Future Work}
\label{sec:limits}

\subsection{Limitations}
\label{sec:limits:lim}

\todo{Be honest about each:

\textbf{Single base model.} Llama-3-8B-Instruct only. Larger and smaller
models may exhibit different sample-efficiency curves; the ranking may not
generalise. Mitigation: report findings as conditional on the base model
and flag this as the most important external-validity threat.

\textbf{Two domains, both well-represented in pretraining.} GoodReads and
Netflix are pop culture, and Llama-3 has demonstrably strong priors on
both. Domains where the LLM has weaker priors -- niche literature,
specialised B2B software, non-English content -- are likely to behave
differently. The familiarity filter (\cref{sec:method:data}) makes this
explicit but doesn't fix it.

\textbf{Preferences derived from explicit ratings.} Real-world preference
data are noisier, sparser, and shaped by survivorship and presentation
effects. The findings here represent an upper bound on what these methods
can achieve.

\textbf{Static, offline setting.} No online adapter updates, no continual
learning, no measurement of how preferences drift over time. All of these
matter operationally and none are addressed here.

\textbf{Compute budget.} Two seeds is enough for direction, not enough
for tight confidence intervals on smaller effect sizes.}

\subsection{Future Work}
\label{sec:limits:future}

\todo{Forward-looking but disciplined:

\textbf{Scale.} Repeat the comparison at $\sim$70B parameters and at smaller
models ($\leq 3$B) deployable on-device. The sample-efficiency curves are
the quantity most likely to change.

\textbf{Domain stress test.} Run the same protocol on a domain where the
base model has thin priors. The familiarity filter would have to be
loosened or replaced.

\textbf{Online and continual alignment.} How does the adapter update as the
user provides more signal over time? Streaming preference learning is a
natural extension of this work.

\textbf{Segment-level adapters.} Per-user adapters do not scale to
millions of users. Clustering users and training one adapter per cluster
is the obvious operationalisation; the question is how much
personalisation quality is lost in the aggregation.

\textbf{Multi-objective alignment.} Real products optimise for preference
fit \emph{and} diversity, safety, and engagement quality. The losses
studied here are single-objective by construction.

\textbf{Cross-domain transfer.} Does aligning on books transfer to films?
The shared latent ``taste'' axis is an open empirical question.}
